% 【基础正文】所属：第8章（由 chapters/revised/08-digital-ai.tex \input 引入）
\minitopic{数字工具改变了什么}搜索引擎和生成式系统没有取消传统认识问题，而是重排了证据链。用户看到的不是世界本身，而是经过抓取、排序、训练、过滤和界面呈现的结果。来源可见性下降，获取速度上升；验证成本往往从“找到内容”转移到“确认内容如何产生”。 需要区分四层：内容是否为真；系统在相关任务上是否可靠；用户是否有理由信任这次输出；组织是否有记录、复核和纠错机制。总体准确率不能保证个案，个案错误也不能证明系统毫无用途。 生成系统可能产生语法连贯却无来源支持的陈述\parencite{bender2021}。NIST把这种风险称为\term{虚构}{confabulation}，即自信生成错误内容，并把它与信息完整性和人机配置风险分别管理\parencite{nist2024}。术语意在描述输出风险，不应拟人化为系统像人一样故意撒谎。

\minitopic{排序、过滤与认识环境}推荐系统选择何者可见，会影响主体拥有的证据集合。个性化可能提高相关性，也可能让已有偏好获得过量确认。所谓“信息茧房”不能只凭不同用户看到不同内容来断定；需要检验来源多样性、跨群体接触、用户主动选择和平台机制。 假信息传播不只源于政治偏见。分析性思考不足、重复暴露、社会认同和情绪唤起都可能提高接受度\parencite{pennycook2019}。治理若只删除错误陈述，可能忽略传播网络如何奖励速度、确定语气和群体忠诚\parencite{oconnor2019}。

\minitopic{算法偏差与责任分配}面部分析系统在不同交叉群体上的误差差异说明，平均性能会遮蔽分布性失败\parencite{buolamwini2018}。认识评价因此至少要问：训练和测试群体是否代表应用对象？错误由谁承担？受影响者能否申诉？模型卡、数据说明和分组评估把部分隐含条件变成可审查记录\parencite{mitchell2019}。 责任不能简单推给“算法”或最终用户。开发者、部署机构、采购者和专业使用者拥有不同控制能力。风险越高，越需要独立验证、来源保留和人工复核；人工在环也不是万能答案，如果人没有时间、权限或专业能力，点击确认只会制造责任幻象\parencite{floridi2019,raji2020}。

\minitopic{机器能够知道吗}若知识只要求可靠地产生真表征，某些系统可能接近功能性知识。若知识要求信念、理解、理由响应、认识责任或主体性，结论会更保守。争论不能仅凭系统说“我知道”决定，因为第一人称输出可能只是语言行为模拟。 更有实践价值的问题是：人机组合能否形成可靠的分布式认识主体？一个系统可检索，人类可提出目标和审查意义，机构可保存来源并纠错。组合表现取决于接口是否展示不确定性、是否允许追溯、用户是否被训练识别失败模式。

\begin{argumentbox}
最低核验流程：把输出拆成可核实主张；追到原始或权威来源；核对题名、作者、日期、页码和DOI；用独立来源复核高风险事实；记录哪些内容由工具生成、哪些由人确认；无法核实时明确标注不确定，而不是用流畅转述填补空白。
\end{argumentbox}

\begin{comparison}
“知之为知之，不知为不知”在数字环境中可转化为校准要求\parencite[2.17]{analects2003}：系统和用户都应使确信程度与证据相称。这里不是说古典文本预见了AI，而是用其认识谦逊规范检验界面设计：系统能否表达不知道，用户能否保留暂停判断的空间？
\end{comparison}

\minitopic{检索与生成不能混为一谈}检索系统返回已有文档的位置，生成系统根据训练和上下文产生新的符号序列；带检索的生成系统则组合二者。看到脚注或链接不能保证句子由链接支持。用户应逐项检查来源是否存在、是否真的蕴涵主张、版本与日期是否匹配。 系统输出的不确定性也有不同含义：模型对下一个词的分布、某项事实为真的概率、系统在任务类型上的历史准确率，并不是同一数值。界面显示“置信度”时必须说明测量对象，否则精确数字会制造伪校准。

\minitopic{自动化偏差与反向偏差}\term{自动化偏差}{automation bias}使人过度接受机器建议；反向的算法厌恶则可能因一次显眼错误放弃总体优于人的工具。合理态度不是总信或总不信，而是按任务、环境与错误类型校准依赖。系统在人类最易犯错的案例上是否也犯同样错误，尤其重要。 人机组合若要优于双方，需要互补信息和真正的否决能力。让疲惫操作员在数秒内确认数百条建议，不构成有意义监督。记录覆写理由、抽样复核和比较人机分歧案例，比简单声明“最终由人负责”更有效。

\minitopic{案例工作坊：AI文献综述}研究者让系统生成十篇核心文献。第一步只把结果当检索候选；第二步在期刊、出版社或权威数据库核对元数据；第三步阅读摘要和正文，确认文献与论点关系；第四步记录系统遗漏和虚构；第五步由研究问题而非生成顺序组织综述。 若系统准确列出九篇而虚构一篇，不能把整份列表直接引用，也不必否认它帮助发现线索。认识地位属于经过核验的人机流程，而不是原始输出。研究者仍对最终断言和引证负责。
